Theorem 2 (Connection to the Shapley value, proved by
[29]). Let φ(i) denote the Shapley value [56] of an input
variablei. Then, the Shapley value φ(i) can be explained
as the result of uniformly assigning causal effects to each
involving variablei, i.e.,φ(i) =∑
S⊆N\{i}
1
|S|+1wS∪{i}.
The Shapley value [56] was first proposed in game theory
and has been used by previous studies [45] to estimate attri-
butions of input variables in the DNN. The Shapley value
satisfies four satisfactory axioms and is widely considered as
a relatively fair estimation of attributions. Theorem 2 proves
that the Shapley value can be considered as a re-allocation
of causal effects to input variables.
In Appendix D.2, we further prove that the Shapley in-
teraction index [26] and the Shapley Taylor interaction in-
dex [66] can also be understood as the assignment of causal
effectswS to different coalitions.
3.2. Discovering and boosting the conciseness of the
causal graph
Remark 1. Given a DNNv(·) and an input sample x withn
variables, we can find a small set of causal patternsΩ subject
to|Ω|≪ 2n, such that the DNN’s output can be approximated
by the causal graph’s output,i.e.∀S⊆N , Y(xS)≈v(xS).
• Discovering the conciseness. We have discovered that
lots of DNNs with various architectures trained for different
tasks can all be explained using sparse causal patterns. Al-
though Theorem 1 indicates that the causal graph needs to
encode 2n causal patterns to precisely fit the DNN’s output
on all the 2n masked samples, Remark 1 shows a common
phenomenon that the causal effects wS extracted from the
DNN are usually very sparse. To this end, we trained various
DNNs for different tasks, and Fig. 2 shows the strength of
causal effects|wS| in descending order for various DNNs.
We found that most causal patterns had little influence on
the output with negligible values|wS|≈ 0, and they were
termed noisy causal patterns . Only a few causal patterns
had considerable effects|wS|, and they were termed salient
causal patterns . Furthermore, we also conducted experi-
ments in Section 4.2, and Figs. 3, 4, and 6 show that we
could use a small number of causal patterns (empirically 10
to 100 causal patterns for most DNNs) in Ω to approximate
the DNN’s output, as stated in Remark 1.
• Boosting the conciseness. Inspired by Remark 1, we
aim to learn a more concise causal graph. To this end, we pro-
pose the following objective of learning faithful and sparse
causal effectswS.
minw,Ω unfaith(wΩ) s.t.|Ω|≤M
⇔ minw,Ω unfaith(wΩ) s.t.‖wΩ‖0≤M,
unfaith(wΩ) =
∑
S⊆N
[
v(xS)−YwΩ(xS)
]2
(5)
where wΩ
def
= [w′
S1,w′
S2,...,w′
S2n ]. IfS∈ Ω, then w′
S =wS;
otherwise,w′
S =0. The L0-norm‖wΩ‖0 refers to the number
of non-zero elements in wΩ, thereby‖wΩ‖0 =|Ω|. In this
way, the above objective function enables people to use a
small number of causal patterns to explain the DNN.
However, direct optimization of Eq. (5) is difficult. There-
fore, we propose several techniques to learn sparse causal
effects based on Eq. (5) to faithfully mimic the DNN’s
outputs on numerous masked samples. The following para-
graphs will introduce how to relax the Harsanyi dividend in
